\section{The Five Functions of an AI-Enabled Workflow}

Watch what any AI tool actually does with a request, and it's typically some blend of five functions. This is a proposed functional lens for evaluating model-mediated applications, useful because it isolates characteristic failure modes --- not a complete or universally applicable taxonomy. Some products do no retrieval, no verification, and take no action at all; others contain training pipelines, memory systems, ranking algorithms, or deterministic computation that don't map neatly onto any one of the five. Real systems are graphs, with parallel paths and blended steps, not a tidy pipeline. Treat what follows as a lens to apply where it fits, not a checklist every system must satisfy.

\textbf{1. Context assembly.} Something decides what information goes in front of the model --- retrieval, memory, prompt construction. This fails \emph{quietly}: missing the one relevant clause produces a confident answer on an incomplete picture, with no visible sign anything went wrong. Much of this step is itself statistical (search, embeddings, reranking), so "gathered context" is a weaker claim than "gathered the \emph{right} context."

\textbf{2. Generation.} The model produces an answer. Generation is the \textbf{distinctive probabilistic component} in these systems, and its errors are frequently what motivate the other four functions to exist as controls --- though for many products, retrieval, deterministic calculation, rule application, or state management are the primary functions, with generation playing a supporting role, or vice versa; it is not correct to treat generation as the one irreducible function every product exists to serve. Its defining property remains real regardless: \textbf{fluency is not a reliable indicator of correctness.} The same polished, confident style accompanies a correct answer, a subtly wrong one, and an outright fabrication.

\textbf{3. Verification.} Something checks the generated answer. \textbf{Governing constraint, stated precisely: verification only catches the kinds of error its checker is built to see.} Four things get blurred under one word: \emph{provenance} (where material came from), \emph{validation} (whether a defined check passed), \emph{evaluation} (measured performance over a representative test set), and \emph{human approval} (who accepts responsibility). A fifth belongs alongside for anything running over time: \emph{monitoring} --- whether performance stays within tested bounds as the model, documents, and usage pattern drift.

For domains involving cited authority, verification splits further: \emph{source grounding} (faithful to documents actually given) versus \emph{authority grounding} (is the cited authority still valid). A tool can pass the former and fail the latter completely.

\textbf{4. Control.} Something decides what happens next --- accept, retry, escalate, act. For consequential deployments, the control step's interface is where responsibility actually lives in practice: a well-designed layer exposes uncertainty to a human reviewer rather than presenting a seamless block of polished text that invites rubber-stamping.

\textbf{5. Effect (action).} The system changes state outside itself --- sends, files, commits. The only function with consequences outside the tool, and correspondingly the one deserving the most caution. A reviewed wrong answer is a mistake caught; an acted-upon wrong answer is a mistake shipped.

\textbf{On generation nesting inside other functions.} A tool that uses a small model to select documents has a generation step hidden inside what looked like mechanical retrieval; a tool that uses a model to grade output has one hidden inside verification. The literature on whether a checker built from the same kind of model as the writer shares its blind spots is genuinely mixed rather than uniformly negative --- some work finds language models struggle to reliably self-correct their own reasoning without external signal, while other work finds self-critique and model-based verification \emph{do} improve results, particularly when the verifier is given external evidence or an independently checkable condition to work from rather than being asked to grade purely on its own judgment. The defensible framing: \textbf{correlated failure between generator and checker is a risk to test for, specific to a given setup, not an invariant to assume.} The dimensions worth checking independence across are the model provider, the specific model family, the prompt, whether external evidence is supplied to the checker, and the type of verifier used (a second LLM call versus a genuinely different mechanism, such as a rule engine or a deterministic tool).

\textbf{On "reasoning."} Visible chain-of-thought does often involve generated intermediate tokens fed back as context, but current reasoning-oriented training also involves reinforcement learning on outcomes, allocation of additional test-time computation, search over candidate solutions, and tool use --- reasoning is not reducible to a single visible application-level loop. \textbf{Iterative chain-of-thought or reflection loops} can genuinely improve performance, but don't by themselves establish reliability, and can --- without an independent intermediate check --- entrench an early mistake with increasing apparent confidence. Additional inference-time computation frequently does improve measured accuracy on a range of tasks, but \textbf{its benefit has to be measured for the specific task rather than assumed, and it does not itself constitute independent verification} --- a model can reason at length toward a wrong answer with just as much apparent confidence as toward a right one.

\textbf{The diagnostic signature of a dangerous configuration:} heavy dependence on generation, paired with verification too weak to cover the risks that actually matter. This demos beautifully and fails silently in production. When a product pitch is strong on the functions easy to show off (context assembly, action) and quiet on the functions hard to build (verification, control), that silence is itself diagnostic information.

\textbf{Silence as evidence, not verdict.} When a claim about a safeguard isn't made, the honest reading is "not established" --- neither "confirmed absent" nor "safe to assume present." Test it, and weight what you find by how it was established: independently verified outranks asserted in a product deck, which outranks merely not contradicted.

\textbf{The five functions are a lens, not the whole assessment.} Whether an organization can actually prevent, detect, and respond to failures depends on cross-cutting controls entirely outside this view: identity and access management, security architecture, contracts, data lifecycle policy, incident response readiness, change management. A tool can be architecturally excellent on all five functions and still unsafe to adopt with no incident plan and no exit route for the organization's data.
